% Figure: a Shewhart control chart for a process mean. The centre line
% is the target; the upper and lower control limits sit at +/- three
% sigma of the sample-mean sampling distribution. Twenty subgroup means
% are plotted in sequence; most sit inside the limits, one point pierces
% the upper limit (an out-of-control signal), and a run of points on one
% side near the end illustrates the run rule that flags a drift before a
% single point ever leaves the band. Coordinates are precomputed.
\begin{figure}[htbp]
\centering
\begin{tikzpicture}
\begin{axis}[
    width=13cm, height=6.8cm,
    xmin=0, xmax=21,
    ymin=94, ymax=106,
    axis lines=left,
    xlabel={Subgroup number},
    ylabel={Subgroup mean $\bar{x}$ (\si{\gram})},
    every axis plot/.append style={mark=none},
    xtick={2,4,6,8,10,12,14,16,18,20},
    ytick={95,98,100,102,105},
]
% Control limits and centre line.
\addplot[engred, very thick] coordinates {(0,105) (21,105)};
\node[anchor=south east, font=\footnotesize, engred] at (axis cs:21,105) {UCL $=105$};
\addplot[engred, very thick] coordinates {(0,95) (21,95)};
\node[anchor=north east, font=\footnotesize, engred] at (axis cs:21,95) {LCL $=95$};
\addplot[engblack, thick] coordinates {(0,100) (21,100)};
\node[anchor=south east, font=\footnotesize, engblack] at (axis cs:21,100) {centre $=100$};
% Subgroup means: in-control run, one out-of-control point, then drift.
\addplot[engblue, thick, mark=*, mark size=1.6pt] coordinates {
 (1,100.4) (2,99.1) (3,101.2) (4,98.6) (5,100.8) (6,99.4)
 (7,101.6) (8,98.2) (9,100.1) (10,105.8) (11,99.7) (12,100.5)
 (13,101.3) (14,101.8) (15,102.1) (16,102.4) (17,102.9) (18,103.2)
 (19,103.6) (20,104.1)};
% Highlight the out-of-control point.
\addplot[only marks, engorange, mark=*, mark size=2.6pt] coordinates {(10,105.8)};
\node[anchor=south, font=\footnotesize, engorange] at (axis cs:10,105.9) {out of control};
\node[anchor=north west, font=\footnotesize, engamber] at (axis cs:14,102.0) {upward run (drift)};
\end{axis}
\end{tikzpicture}
\caption[Shewhart control chart with two out-of-control signals]{A
Shewhart chart plots subgroup means in production order against a
centre line at the target and control limits at three standard errors
of the subgroup mean. A single point outside the limits, subgroup ten
here, is the classical out-of-control signal and triggers an
investigation. The run of eight rising points near the end pierces no
limit yet, but the supplementary run rules flag it as a drift, catching
a slow shift of the process mean before any single point escapes the
band. The chart is hypothesis testing turned into a live procedure:
each plotted point is a test of the null that the process is on target,
and the three-sigma limits set a per-point false-alarm rate near
$0.0027$.}
\label{fig:vol02:ch14:control-chart}
\end{figure}
